\subsection{Monodromies on the Moduli Space of the $SU(2)$ Theory:}

We have discussed the possibility of the existence of monodromies on
the moduli space and pointed out that understanding the monodromy
structure may help us solve the theory, {\it i.e.,} to determine the
exact non-perturbative prepotential for the effective low-energy
theory. This is equivalent to determining the functions $a_D(u)$ and
$a(u)$. In this subsection, we set out to identify the singularities
on the moduli space ${\cal M}$ and calculate the associalted
monodromies.  

\noindent\underline{Monodromy at large $u$:} 
At large $|a|$, the theory is asymptotically free and
$u=\frac{1}{2}a^2$. To a good approximation, the prepotential is given
by the one-loop formula ${\cal F}(a)=(i/2\pi)a^2\,{\rm ln}\,
(a^2/\Lambda^2)$, from which we obtain 
$$
a_D=\frac{\partial{\cal F}}{\partial a}=
\frac{2ia}{\pi}\,{\rm ln}\,(\frac{a}{\Lambda})+\frac{ia}{\pi}\,.
$$
If we make a close loop on the $u$-plane around $u=0$, we get 
${\rm ln}\,u \rightarrow {\rm ln}\,u+2\pi i$, and ${\rm ln}\,a
\rightarrow {\rm ln}\,a+i\pi$, hence, 
$$
\begin{array}{rcl}
a_D & \rightarrow & - a_D + 2a \,,\\
a & \rightarrow & - a \,.
\end{array}
$$
This is implemented by the monodromy matrix
\begin{eqnarray}
M_{\infty}&=& PT^{-2}=\left(\begin{array}{cc}
-1 & 2 \\ 0 & -1 \end{array} \right )\,,
\label{mondone}
\end{eqnarray}
acting on $(a_D, a)^T$.  Here, $P$ is the negative of the identity
element in $SL(2,{\rm Z})$ and $T$ is as defined in (\ref{ST}). Under
the action of this monodromy, the magnetic and electric quantum
numbers of BPS states change as $(n_m,n_e)\rightarrow
(-n_m,-n_e-2n_m)$ so that the mass formula is unchanged.

\noindent\underline{Monodromies at finite $u$:}
The monodromy at $u=\infty$ implies that there exist other monodromies
somewhere else on the $u$-plane. If these monodromies commute with
$M_\infty$, then $a^2$ is a good global coordinate. However, this
cannot be the case as then the positivity of the kinetic energy is
violated. To get a non-abelian monodromy group, we need at least two
singularities on the $u$-plane, and at finite $u$, with non-trivial
monodromies around them. These singularities will be related by the
broken discrete symmetry $u\rightarrow -u$. We make this minimal
assumption on the number of extra singularities. Then a loop enclosing
both these singularities should reproduce the monodromy $M_\infty$

What is the origin of these singularities on the moduli space? To
obtain the Wilsonian low-energy effective action, we have integrated
out all massive states in the theory. The massive particle loops then
give rise to a non-trivial metric on the moduli space. The values of
the masses, however, depend on the modulus $u$, and it may so happen
that for certain values of $u$, some of the masses become zero. Then,
at these points, we end up integrating out massless states and thus
create singularities on the moduli space. The nature of a singularity,
{\it i.e.,} the monodromy associated with it, depends on the
properties of the particle which becomes massless at the singularity.
Since, for finite $u$ (the non-perturbative regime), the mass spectrum
as a function of $u$ is not known, the singularities cannot be found
in a straightforward way. The way to proceed is to assume that some 
generic states become massless at certain values of $u$ (say $u=1$ and
$u=-1$) and find the corresponding monodromies (say $M_1$ and
$M_{-1}$). The massless states are then specified by the condition  
$M_1 M_{-1}=M_\infty$.

Naively, one may expect that the massive gauge boson multiplets
contribute to the singularity when they become massless. Classically,
this happens at the point $u=0$, which, in quantum theory, may
get shifted to a non-zero $u$. However, as argued in \cite{SW-I}, a
spin-1 multiplet becoming massless does not lead to a consistent
picture. We will not repeat this argument here.  

The only other massive states in the theory with spin $\le 1/2$ are
monopoles and dyons, which, due to the spin condition, belong to short
$N=2$ multiplets and, therefore, are BPS states. These states are
described by hypermultiplets. Seiberg and Witten suggested that the
singularities arise when some of these non-perturbative states become
massless.  The problem now is to calculate the associated
monodromies. Note that the hypermultiplets for monopoles and dyons
cannot be coupled to the fundamental fields of our theory in a local
way. However, in the subsection on duality it was seen that it is
possible to go to a dual description of the theory in which some dual
gauge fields couple to the monopoles or dyons in exactly the same way
that the usual gauge fields couple to a particle of unit electric
charge. Thus we only have to calculate the monodromy when a massive
electrically charged hypermultiplet becomes massless and then, using
the duality transformation, find the monodromy for a generic monopole
or dyon. Let us consider a dual description of the theory in which a
certain monopole or dyon appears as an elementary state, and let us
label this description by a letter, say $q$. Near the point where this
state is massless, all massive fields can be integrated out and the
theory is essentially a $U(1)$ theory coupled to a hypermultiplet. If
we denote the {\it vev} of the scalar field in this description of the
theory by $a(q)$, then the mass of a BPS state of unit electric charge
goes to zero when $a(q)=0$ at some $u=u_q$. Thus, near this point,
$a(q)$ is a good local coordinate and can be expanded as $a(q)\approx
c_q\,(u-u_q)$. Moreover, near this point, the one-loop $U(1)$
$\beta$-function implies (see the next subsection):
$$
\tau (a(q))=-\frac{i}{\pi}\,{\rm ln}\,\frac{a(q)}{\Lambda}\,,
$$
from which we obtain
$$
a_D(q)=-\frac{i}{\pi}\,a(q)\,{\rm ln}\,\frac{a(q)}{\Lambda}+
\frac{i}{\pi} \,.
$$
Moving on a closed loop around $u_q$ so that $(u-u_q)\rightarrow
e^{2\pi i} (u-u_q)$, we get the monodromy
\begin{equation}
\begin{array}{rcl}
a_D(q)& \rightarrow & a_D(q) + 2 a(q)\,,\\
a(q)  & \rightarrow & a(q)\,.
\end{array}
\label{mondq}
\end{equation}

Let us now calculate the monodromy when a $(n_m,n_e)$ dyon becomes
massless (the dyon is stable or marginally stable if $n_m$ and $n_e$
are coprime). The first step is to find a dual description of the
theory in which this dyon appears as an elementary stat of charge
$(0,1)$. Under a generic $SL(2,{\rm Z})$ transformation we get a 
$\left(n_m(q),n_e(q)\right)$ dyon with
\begin{equation}
\left(\begin{array}{c}a_D(q)\\a(q)\end{array}\right)=
\left(\begin{array}{c}\alpha a_D + \beta a\\
                      \gamma a_D + \delta a \end{array}\right)\,,
\qquad
\left(\begin{array}{c}n_m(q)\\n_e(q)\end{array}\right)=
\left(\begin{array}{c}n_m\delta - n_e\gamma \\
                    -n_m\beta + n_e\alpha \end{array}\right)\,,
\label{asns}
\end{equation}
so that $Z=n_m a_D + n_e a$ is invariant. In the above, $\alpha\delta
- \beta\gamma =1$. Now we choose the parameters such that $n_m(q) =0,
n_e(q)=1$. With this choice, $(a_D(q), a(q))$ become the variables in
terms of which the dyon couples to the $SL(2,{\rm Z})$ transformed
gauge field in the same way that the unit electric charge couples to
usual gauge fields. In particular, when the dyon becomes massless at
some point on the moduli space, the associated monodromy, in this
description, is given by (\ref{mondq}). Inverting the first equation
in (\ref{asns}), we get $a_D=-\beta a(q) + n_e a_D(q)$ and $a=\alpha
a(q) - n_m a_D (q)$. Thus we can easily find the action of the
monodromy on the original variables as 
\begin{equation}
\left(\begin{array}{c}a_D<本地路径>)\rightarrow 
\left(\begin{array}{cc}1+2n_e n_m & 2n_e^2\\
-2n_m^2 & 1-2n_e n_m \end{array}\right)\,
\left(\begin{array}{c}a_D<本地路径>)\,.
\label{mondn}
\end{equation}
Denoting the monodromy matrix as $M(n_m,n_e)$, we note that ${\rm
  Tr}\,M(n_m,n_e)=2$. Thus the monodromy always belongs to the parabolic
subgroup of $SL(2,{\rm Z})$. 

Now we calculate the monodromies at $u=\pm 1$. Let us assume that a
$(m,n)$ dyon becomes massless at $u=1$ and a $(m',n')$ dyon becomes
massless at $u=-1$. The the associated monodromies should satisfy
\begin{equation}
M_1(m,n)\,M_{-1}(m',n') = M (\infty)\,.
\label{mondeq}
\end{equation}
Using (\ref{mondone}) and (\ref{mondn}), this can be written as 
$$
\left ( \begin{array}{cccc} 1 + 2mn & 2n^{2} \\ -2m^{2} & 1-2mn
\end{array} \right ) =
\left( \begin{array}{cccc} -1 & 2 \\ 0 & -1 \end{array} \right )
\left ( \begin{array}{cccc} 1-2m'n' & -2n'^{2} \\
2m'^{2} & 1 + 2m'n' \end{array} \right )\,,
$$
leading to the equations
$$
\begin{array}{rcl}
1 + mn &=& m'n' + 2m'^{2} \,,\\
m^{2} &=& m'^{2} \,,\\
n^{2} &=& n'^{2} + 1 + 2m'n'\,, \\
1 - mn &=& -m'n'\,.
\end{array}
$$
These imply that $m=\pm 1$ and $m'=\pm 1$. For each combination of
$(m,m')$, we can easily determine $n'$ in terms of $n$ and get the
following possible sets of solutions
$$
\begin{array}{cccccccc}
(m,n): &(1,n)&,&(-1, n)&,&(-1, n)&,&(1, n)\,,\\
(m',n'): &(1, n-1)&,&(1, -n-1)&,&(-1, n+1)&,&(-1, -n+1)\,. 
\end{array}
$$
There do not seem to be any solutions where $M_{\infty}$ is factorized
into a product of more than two such parabolic $M(m,n)$. Moreover,
note that the solution allows only dyons of unit magnetic charge to
contribute to the monodromy. This is  consistent with the result that,
semiclassically, only these dyons are stable. 

Let us consider another consistency check: In general, under the
action of the monodromy, the quantum numbers of dyons will
change. However, we expect that the particular dyon which becomes
massless and is the source of the singularity should remain invariant
under the monodromy (as it is the properties of this dyon which
determine the monodromy matrix). This can be easily checked. The
eigenvalue equation
\begin{equation}
(q_m, q_{e}) \left( \begin{array}{cc}
1+2mn & 2n^{2} \\ -2m^{2} & 1-2mn \end{array} \right ) =
(q_m, q_{e})\,,
\label{mon-inv}
\end{equation}
leads to $nq_m-mq_e=0$, provided $m$ and $n$ do not vanish
simultaneously. Clearly, $q_m=m, q_n=n$ is a solution of this
equation. If we restrict ourselves to stable dyons, then this is the
only solution. Thus, knowing the monodromy, we can find the dyon which
gives rise to it.   

The simplest solution to the equation (\ref{mondeq}) corresponds to 
$m=m'=1, n=0, n'=-1$. The monodromy matrices are then given by 
\begin{equation}
M_1=\left(\begin{array}{cc}1&0\\-2&1\end{array}\right),\qquad
M_{-1}=\left(\begin{array}{cc}-1&2\\-2&3\end{array}\right)\,.
\label{mondothers}
\end{equation}
A comparison with (\ref{mon-inv}) implies that $M_1$ arises due to a
monopole becoming massless at $u=1$ and $M_{-1}$ arises when a
$(1,-1)$ dyon becomes massless at $u=-1$. At the point where the
monopoles become massless, we have $a_D=0$ and where the $(1,-1)$ dyon
becomes massless, we have $a-a_D=0$ as is evident from $Z=n_m a_D +
n_e a$. Note that the monodromy at infinity, $M_\infty$, shifts the
electric charge by $2$ units. Hence, at the points where condensation
takes place, the electric charge is really defined modulo $2$. We can
conjugate the representation of the fundamental group by
$M^{n}_{\infty}$.

\subsection{$U(1)\; \beta$-function}

To construct the monodromies in the previous subsection, we have used
the $\beta$-function for a $U(1)$ theory interacting with a
hypermultiplet. Let us look at this in some more detail. If we have
Weyl fermions with charge $Q_{f}$ and complex scalars with charge
$Q_{s}$, then their contribution to the $U(1) \beta$-functions is
$$
\beta(g)\equiv\mu \frac{d}{d \mu}g = \frac{g^{3}}{16\pi^{2}} \,
\left( \sum_{f} \; \frac{2}{3} Q^{2}_{f} +
\sum_{s} \; \frac{1}{3} Q^{2}_{s} \right )\,.
$$
If we denote the coefficient of $g^3$ by $b$ and define $\alpha= 
g^2/4\pi$, then the above equation can be rewritten as
$$
\mu \frac{d}{d\mu} \left( \frac{1}{\alpha} \right) = - 8\pi b\,.
$$ 
Consider hypermultiplet which is a reduced multiplet of $N = 2$ with
spin $\leq 1/2$. In $N=1$ language, this is described by chiral
superfields $M,\widetilde M$ and contains two Weyl fermions and two
complex scalars, all with the same charge $Q$. Hence we get
$$
b = \frac{1}{16\pi^{2}} \; Q^{2} \left(2 \cdot \frac{2}{3} +
2 \cdot \frac{1}{3} \right) = \frac{1}{8\pi^{2}}\; Q^{2}\,.
$$
Now remember that, using the anomaly, we have set the
$\theta$-parameter to zero by a chiral rotation of the fermions. As a
result we have $\tau=i/\alpha$, so that, 
$$
\mu \frac{d\tau}{d\mu} = - \frac{i}{\pi} Q^{2}\,.
$$
Identifying $\mu$ with the natural scale of the theory which is $a$
and setting $Q=1$, we obtain
$$
\tau \simeq - \frac{i}{\pi}\,{\rm ln}\,\frac{a}{\Lambda}\,.
$$
This is the expression which was used in the determination of the
monodromies at finite $u$. If we are interested in the contribution
from a monopole multiplet, then we can perform the above calculation
in terms of the dual variables. The answer then becomes $\tau_D\simeq
-(i/\pi)\, {\rm ln}\,(a_D/\Lambda)$.   

\subsection{Monopole Condensation and Confinement}

In this subsection we describe how confinement in $N=1$ gauge theory
can be understood in terms our macroscopic $N=2$ theory. The $N = 2$
vector superfield ${\cal A}$ can be decomposed, in the $N=1$ formalism,
into a vector superfield $W_{\alpha}$ and a chiral superfield $\Phi$. To
break $N = 2$ down to $N = 1$, one can add a superpotential $W=m\;
{\rm Tr}\Phi^2$ to the action. This gives a bare mass to the $\Phi$
multiplet and, therefore, the low-energy theory is a pure Abelian
gauge theroy. This low-energy theory has ${\bf Z}_{4}$ chiral
symmetry and is believed to have a mass gap with confinement 
and spontaneous breaking of ${\bf Z}_4$ to ${\bf Z}_2$. Seiberg and
Witten gave a macroscopic description of this phenomenon based on the
$N=2$ picture as will be described below.

In $N=2$ theory, the massless spectrum in the semiclassical limit
contains only the Abelian vector multiplet ${\cal A}$.  Let us analyze
the effect of turning on a mass $m$ for the scalar multiplet
$\Phi$. In the low-energy theory, ${\rm Tr}\,\Phi^2$ is represented by
a chiral superfield $U$. Its scalar component is $u =\langle{\rm
Tr}\phi^2\rangle$, which is a holomorphic function on the moduli
space. For small $m$, we can simply add $W_{eff} = mU$ to the 
low-energy Lagrangian. This presumably removes the vacuum degeneracy
and gives a mass to the scalar multiplet. To make the Abelian gauge
field also massive (so that the theory has a mass gap), we need either
(i) extra light gauge fields giving rise to a strongly coupled
non-Abelian theory, or (ii) light charged fields giving rise to a
Higgs mechanism.  Thus, in either case, somewhere on ${\cal M}$ extra
massless states must appear. As indicated before, one cannot get extra
light gauge fields. Therefore, we consider option (ii) with the light
charged fields being monopoles or dyons. Near the point with massless
monopoles, we go to a dual description of the theory and use the $N=1$
chiral superfields $M, \widetilde M$ to describe the monopole
hypermultiplet. In this dual description, the full $N=1$
superpotential then becomes
$$
\hat{W} = \sqrt{2} A_{D} M \widetilde M + m \; U(A_{D})\,.
$$
The low-energy vacuum structure is easy to analyze. Vacua correspond
to solutions of
$$
d \hat{W} = 0\,,
$$
satisfying $|M| = |\widetilde M|$, so that the $D$-term vanishes. For $m
=0$, $M=\widetilde M =0$ and $a_{D}$ is arbitrary. Thus we recover the
$N=2$ moduli space. If $m \not= 0$, then 
\begin{eqnarray}
&{\displaystyle
\sqrt{2} M \widetilde M + m \; \frac{du}{dA_{D}} = 0} &\,,\nonumber \\
& a_{D}M = a_{D}\widetilde M = 0 &\,.\nonumber 
\end{eqnarray}
Assuming that $du/da_{D} \not= 0$, we get $M, \widetilde M\not= 0$, so
that $a_D=0$ and $M=\widetilde M=(-m u'(0)/\sqrt{2})^{1/2}$. Since $M$
is charged, its vacuum expectation value generates a mass for the
gauge field through the Higgs mechanism. Thus, by a simple analysis,
we reproduce the expected mass gap of the microscopic theory. To
understand charge confinement, note that, since the hypermultiplet
$M\widetilde M$ describes monopoles, we have a magnetic Higgs
mechanism. Thus, $M\ne 0$ means that massless magnetic monopoles
condense in vacuum. This gives rise to the confinement of electric
charges by the dual ({\it i.e.,} magnetic) Meissner effect.  

\subsection{The Solution of the Model}

In this subsection, we will identify the moduli space of the $N=2$
theories with the moduli space of genus 1 Riemann surfaces and then
use this identification to calculate $a_D(u)$ and $a(u)$. 

Let us first summarize what we have learnt about the structure of the
moduli space: The moduli space ${\cal M}$ is the $u$-plane with
singularities at $1, -1, \infty$ and a ${\bf Z}_{2}$ symmetry acting
as $u \rightarrow -u$ (Fig. 3). Over this punctured plane there is a
flat $SL(2,{\bf Z})$ bundle $V$, which has $(a_D, a)^T$ as a
section. This bundle has monodromies  
\begin{equation}
M_{1} = \left( \begin{array}{cc}
1 & 0 \\ -2 & 1 \end{array} \right )\,, \quad
M_{-1} = \left( \begin{array}{cc}
-1 & 2 \\ -2 & 3 \end{array} \right )\,, \quad
M_{\infty} = \left( \begin{array}{cc}
-1 & 2 \\ 0 & -1 \end{array} \right )\,,
\label{mondall}
\end{equation}
around the singularities. 
% ==== FIGURE_SLOT_3 ====

To be more precise, the quantities $(a_D(u), a(u))$ form a holomorphic
section of the bundle $V\otimes {\bf C}$, and have the following
asymptotic behaviour
\begin{eqnarray}
u\approx \infty\,:&\qquad& \left\{\, \begin{array}{l}
a \cong \sqrt{2u}\,  \\
a_D \approx i\;\frac{\sqrt{2u}}{\pi}\;{\rm ln}\; u
\end{array} \right. \,, \nonumber\\
u=1\,:&\qquad& \left\{\, \begin{array}{l}
a_D \approx \,c_0 (u-1) \\
a \approx a_0 + \frac{i}{\pi} \; a_{D}{\rm ln}\; a_{D}
\end{array} \right. \,.\nonumber
\end{eqnarray}
where, $a_0$ and $c_0$ are constants. For $u=-1$ we get a behaviour
similar to $u=1$ but with $a_D$ replaced by $a-a_D$. The metric on the
moduli space is $ds^2={\rm Im}(\tau)\,|da|^2$ with 
\begin{equation}
\tau =\frac{da_D/du}{da / du}\,.
\label{tau-metric}
\end{equation}
To insure positivity of kinetic energy, ${\rm Im}(\tau)$ should be
positive definite. The monodromies generate a subgroup $\Gamma(2)$ of
$SL(2,{\rm Z})$ and, in fact, the $u$-plane with its singularities is
the quotient of the upper half plane $H$ by $\Gamma(2)$. This quotient
has three cusps corresponding to the three singularites.

The space $H/\Gamma(2)$, which is the $u$-plane, also parametrizes the
family of curves $E_u$ described by the equation 
\begin{equation}
y^2 = (x-1)(x+1)(x-u)\,, 
\label{cubic}
\end{equation}
where $x$ is a complex variable. First, note that this equation is
invariant under the transformations  
$$
w:\{u\rightarrow -u\,,x\rightarrow -x\,,y\rightarrow \pm iy \}\,,
$$
that generate a ${\bf Z}_{4}$ symmetry, out of which only a ${\bf
Z}_{2}$ subgroup acts on $u$. This is the same as the symmetry
structure on ${\cal M}$. Now, let us describe the curve which the
above equation represents. The curve basically is the $x$-space the
topology of which is determined by the requirement that $y$ is a
single valued function. Sine the equation is quadratic in $y$, if we
move along a close loop on the $x$-space around anyone of the three
zeros of $y$, then we get $y\rightarrow -y$. The same is also true for
a loop which contains all the zeros, or equivalently, a loop around
the point $x=\infty$ (this is because there is an odd number of
zeros).  Therefore, if $y$ is to be a single valued function, then the
$x$-space should be a double cover of the complex plane ${\bf C}$ with
the point at infinity added to it. Furthermore, this space should have
four branch points at $x=-1,1,u,\infty$ which are joined pairwise by
two cuts (Fig. 4). To fix attention, consider one branch cut from $-1$
to $1$, and another from $u$ to $\infty$. The two  sheets are joined
along these cuts so that on crossing a cut, we move from one sheet to
the other. It is on this space that $y$ is single valued.

% ==== FIGURE_SLOT_4 ====

The $x$-space so obtained is nothing but a genus one Riemann
surface. This can be easily visualized as follows: On a torus, draw a
circle $c_1$ along the $a$-cycle and translate this along the torus to
get a circle $c_2$ (Fig. 5(a)). Now, squash the circles $c_1$ and
$c_2$ into line segments $l_1$ and $l_2$ (Fig. 5(b)). This divides the
torus into two halves joined along these segments. If we now cut open
both the two halves, the surface we get is the same as the $x$-space
described above with $l_1$ and $l_2$ as the two branch cuts and with
the point at infinity mapped to a point at finite $x$ (Fig. 5(c)).

% ==== FIGURE_SLOT_5 ====

From this, it is clear that a loop on the $x$-plane that goes around
one of the two cuts corresponds to the $a$-cycle of the torus. On the
other hand, a loop which intersects both cuts (in our case, goes
around $1$ and $u$), corresponds to the $b$-cycle on the torus.
Note that when two of the branch points on the $x$-space coincide, a
cycle on the curve $E_u$ shrinks to zero size and the curve becomes
singular. For example, in our case, if $u=\infty$, the $a$-cycle 
shrinks to zero size and when $u=1$, this happens to the $b$-cycle.
Thus the singularities on the $u$-plane are at the points where a
curve in the family $E_u$ develops a vanishing cycle.

To identify $a_D$ and $a$, on the genus one Riemann surface $E_u$, we
pick up two independent one cycles $\gamma_1$ and $\gamma_2$,
normalized such that their intersection number is one: 
$$
\gamma_1 \cdot \gamma_2 = 1\,.
$$
These cycles, which continuously vary with $u$, form a local basis for
the first homology group $V_u=H^1(E_u,{\rm C})$ of $E_u$. A cycle can
be paired with elements $\lambda$ from the first cohomology group 
$$
\gamma \rightarrow  \oint_{\gamma} \lambda\,.
$$
$\lambda$ can be thought of as a meromorphic $(1,0)$-form on $E_u$
with vanishing residue, modulo exact forms. The vanishing of the
residue insures that the pairing is invariant under continuous
deformations of $\gamma$ even across poles of $\lambda$. By virtue of
this pairing, we can also regrad $\lambda$ as an element of $V_u$. For
the one-forms on $E_u$, we can choose a basis   
$$
\lambda_{1} = \frac{dx}{y}\,, \qquad \lambda_{2} = \frac{xdx}{y}\,.
$$
Upto scalar multiplication, $\lambda_1$ is the unique holomorphic
differential on $E_u$ and if we define
$$
b_i=\oint_{\gamma_i}\lambda_1\,,\qquad{\rm for}\quad i = 1,2\,,
$$
then the torus is characterized by a parameter
$$
\tau_u = b_1/b_2\,,\qquad {\rm with} \quad {\rm Im}(\tau_u) > 0\,. 
$$
Let us consider an arbitrary section
$$
\lambda = a_1(u) \lambda_1 + a_2 (u) \lambda_2\,,
$$
of $V_u$ and, for the moment, make the identification
$$
a_D=\oint_{\gamma_1}\lambda\,,\qquad a=\oint_{\gamma_2}\lambda\,.
$$
If $\lambda$ is a form with vanishing residue then, on circling a
singularity, $a_D$ and $a$ transform in the right way, simply
according to how $\gamma_1$ and $\gamma_2$ transform under a subgroup
of $SL(2,Z)$ (this is further explained in subsection 4.6 below). On
the other hand, if $\lambda$ has a pole with a non-vanishing residue,
then it is possible that the integration path may move across this
pole and, as a result, $a_D$ and $a$ may no longer transform under a
pure $SL(2,Z)$. This second possibility is of course not consistent
with the symmetries of the BPS mass formula in the absence of fermions
with non-zero bare masses, as we already discussed in a previous
subsection. Therefore, $\lambda$ should not have a pole with a
non-vanishing residue (In the presence of fermions with non-zero bare
masses, the situation is different and will be discussed at the end of
the next section). The above identification of $a_D$ and $a$ implies
that  
$$
\frac{da_D}{du}=\oint_{\gamma_1}\frac{d\lambda}{du}\,,\qquad
\frac{da}{du}=\oint_{\gamma_2}\frac{d\lambda}{du}\,.
$$
To fix the arbitrariness in $\lambda$, we use the condition 
${\rm Im}\tau > 0$ for the metric on ${\cal M}$ as defined
in (\ref{tau-metric}). First, suppose that 
$$
\frac{d\lambda}{du} = f(u) \lambda_{1} = f(u) \frac{dx}{y}\,.
$$
Then,
$$
\frac{da_{D}}{du} = f(u)\,b_1\,,\qquad \frac{da}{du} = f(u)\,b_2\,,
$$
so that
$$
\tau = \frac{b_1}{b_2} = \tau_u\,. 
$$
Since ${\rm Im}\,\tau_u > 0$, we get ${\rm Im}\,\tau >0$. As argued by
Seiberg and Witten, the converse is also true, so $d\lambda/du$ does
not depend on $\lambda_2$. The function $f(u)$ is fixed by the
asymptotic behaviour of the theory near the singularities on the $u$
plane and is given by $f(u)=-\sqrt{2}/4\pi$. With this, we can obtain
$\lambda$ as
$$
\lambda=\frac{\sqrt{2}}{2\pi}\;\frac{dx\sqrt{x-u}}{\sqrt{x^2-1}}=  
\frac{\sqrt{2}}{2\pi}\;\frac{dx y}{x^2-1}=\frac{\sqrt{2}}{2\pi}\;
\frac{dx}{y} (x-u) \,.
$$
To calculate $a$ and $a_D$, we have to choose a specific basis of 
one-cycles on $E_u$. We identify $\gamma_2$ with the $a$-cycle on the
torus, or equivalently, with a curve which loops around the points
$-1$ and $1$ on the $x$ plane. We can deform this curve so that it lies
entirely along the cut from $-1$ to $1$ and back. Thus, $a(u)$ is
given by 
\begin{equation}
a(u)=\frac{\sqrt{2}}{\pi}\,\int_{-1}^1\,\frac{dx\sqrt{x-u}}{\sqrt{x^2-1}}\,.
\label{au}
\end{equation}
For $\gamma_1$, we choose the curve which loops around the points $1$
and $u$ and get
\begin{equation}
a_D(u)=\frac{\sqrt{2}}{\pi}\,\int_{1}^u\,\frac{dx\sqrt{x-u}}{\sqrt{x^2-1}}\,.
\label{adu}
\end{equation}
It can be checked that with this choice of the one-cycles, $a$ and
$a_D$ have the desired  behaviour near the singularities \cite{SW-I}. 

\noindent\underline{Explicit Formulae for $a(u)$ and $a_D(u)$}:
One can easily find explicit formulae for the $a$ and $a_D$ in terms
complete elliptic integrals $E(k)$ and $K(k)$ by using the integral
representation of hypergeometric functions $F(\alpha,\beta,\gamma;z)$ 
\cite{GR}. The hypergeometric functions are given by
\begin{eqnarray}
F(\alpha,\beta,\gamma;z)&=& \frac{\Gamma(\gamma)}{\Gamma(\beta)
\Gamma(\gamma-\beta)}\int_0^1\,dx\,x^{\beta-1} (1-x)^{\gamma-\beta-1}
(1-zx)^{-\alpha}\nonumber\\ 
&=& \frac{\Gamma(\gamma)}{\Gamma(\alpha)\Gamma(\beta)}\sum_{n\geq 0}
\frac{\Gamma(\alpha +n)\Gamma(\beta +n)}{\Gamma(\gamma +n)}\,
\frac{z^n}{n!}\,.\nonumber
\end{eqnarray}
Comparing this with the integral representation of $a$ in (\ref{au}),
we can easily see that 
$$
a(u)=\sqrt{2(1+u)}\, F\left(-1/2, 1/2, 1; 2/(1+u)\right)\,.
$$
As for $a_D$, let us first make the substitution $x=(u-1)t+1$ in 
(\ref{adu}). This gives
$$
a_D(u)=\frac{i}{\pi} (u-1)\int_0^1\,dt\,t^{-\frac{1}{2}}
(1-t)^{\frac{1}{2}} (1-\frac{1-u}{2} t)^{-\frac{1}{2}}\,. 
$$
Comparing this with the expression for the hypergeometirc functions,
we get
$$
a_D(u)=\frac{i}{2} (u-1) F\left(1/2, 1/2, 2; (1-u)/2\right)\,.
$$

In terms of the hypergeometric functions, the complete elliptic
integrals are given by 
\begin{eqnarray}
K(k) &=& \frac{\pi}{2}\, F\left(1/2, 1/2, 1; k^2\right)\,,\nonumber\\
E(k) &=& \frac{\pi}{2}\, F\left(-1/2, 1/2, 1; k^2\right)\,.\nonumber
\end{eqnarray}
Further, we define $k'^2=1-k^2$ and set 
$$
E'(k)\equiv E(k')\,,\qquad K'(k) \equiv K(k')\,.
$$ 
Now, using the identity
$$
c(1-z)F(a,b,c;z) - c F(a,b-1,c;z) + (c-a)z F(a,b,c+1;z)=0\,,
$$ 
with $c=1, a=b=1/2$ and $z=(1-u)/2$, in the expression for $a_D$, we
get 
\begin{eqnarray}
a(u)   &=& \frac{4}{\pi k} E(k)\,,\nonumber\\
a_D(u) &=& \frac{4}{i\pi} E(i\frac{k'}{k}) + 
\frac{4i}{\pi k^2} K(i\frac{k'}{k})\,.\nonumber
\end{eqnarray}
Here, $k^2=2/(1+u)$ and $(1-u)/2=-k'^2/k^2$. These expressions can be
further simplified if we note that  
$$
F(a, b, c; z)=(1-z)^{-a} F(a, c-b, c; z/(z-1))=
(1-z)^{-b} F(b, c-a, c; z/(z-1))\,.
$$
This implies that $K(ik'/k)=kK(k')=kK'(k)$ and $E(ik'/k)=k^{-1}E(k')
=k^{-1}E'(k)$, so that
\begin{eqnarray}
a(u)  &=& \frac{4}{\pi k}\, E(k)\,,\nonumber\\
a_D(u)&=& \frac{4}{\pi i}\,\frac{E'-K'}{k}\,.\nonumber
\end{eqnarray}
The generalized coupling $\tau$ can be calculated by using
\begin{eqnarray}
\frac{dE}{dk}&=&\frac{E-K}{k}\,,\quad 
\frac{dK}{dk}=\frac{1}{kk'^2}\,(E-k'^2 K)\,,\nonumber\\
\frac{dE'}{dk}&=&-\frac{k}{k'^2}\,(E'-K')\,,\quad 
\frac{dK'}{dk}=-\frac{1}{kk'^2}\,(E'-k^2 K')\,,\nonumber
\end{eqnarray}
and is given by
$$
\tau=\frac{\partial a_D}{\partial a}=\frac{da_D/dk}{da/dk}
=\frac{iK'}{K}\,.
$$
